\subsection{Hypergeometric and Negative Binomial Probability Distributions}
\hrulefill

[Catch up on negative binomial info from last class]
Summary: (to be refined later)
Run n independent bernoulli trials
Decide you require r successes
X = number of trials required to obtain r successes
$X\sim NB(r,p)$
$X = r, r+1, r+2, ...$
$P(X=k) = \binom{k-1}{r-1} p^r (1-p)^{k-r}$
$E(X) = \frac{r}{p}$
$V(X) = \frac{r(1-p)}{p^2}$

\hrulefill

\subsubsection*{Geometric Probability Distribution}
\paragraph*{Definition:}
This is a special case of the negative binomial distribution where $r=1$. Essentially we are looking for the number of trials required 
to obtain the first success. This is denoted:
\[
    X \sim \text{Geometric}(p)
\]
The pmf is given by:
\[
    P(X=k) = p(1-p)^{k-1}, \quad k=1,2,3,...
\]
The mean and variance are given by:
\[
    E(X) = \frac{1}{p}, \qquad V(X) = \frac{1-p}{p^2}
\]

\paragraph*{Review of Geometric Series:}
\[
    1+A+A^2+A^3+...+A^n = \begin{cases}
        \frac{1-A^{n+1}}{1-A}, & A \neq 1\\
        n+1, & A=1 \\
    \end{cases}\frac{1-A^{n+1}}{1-A}
\]

\[
    \sum_{k=0}^{\infty} A^k = \frac{1}{1-A}, \quad |A|<1
\]

\paragraph*{Example:}
You roll a fair die until you see the first six. Define the random variable you are working with in words and state 
the corresponding distribution along with all relevant parameters. (a) How many rolls do you expect to make? (b) What is the 
probability that the first six occurs on the fifth roll? (c) What  is the probability that you will need more than 10 rolls?

Let $X \sim \text{Geo}(\frac{1}{6})$ be the number of rolls until the first six.
\paragraph*{a.}

\[
    E(X) = \frac{1}{p} = \frac{1}{\frac{1}{6}} = 6
\]

\paragraph*{b.}
\begin{align*}
    P(X=5) &= p(1-p)^{5-1} \\
    &= \frac{1}{6} \left(1-\frac{1}{6}\right)^4 \\
    &= \frac{1}{6} \left(\frac{5}{6}\right)^4 \\
    &= \frac{625}{7776} \approx 0.0804
\end{align*}

\paragraph*{c.}
\begin{align*}
    P(X>10) &= 1 - P(X \leq 10) \\
    &= 1 - \sum_{k=1}^{10} P(X=k) \\
    &= 1 - \sum_{k=1}^{10} p(1-p)^{k-1} \\
    &= 1 - p \sum_{k=0}^{9} (1-p)^k \\
    &= 1 - p \left(\frac{1-(1-p)^{10}}{1-(1-p)}\right) = 1 - (1-(1-p)^{10}) \\
    &= (1-p)^{10} = \left(1-\frac{1}{6}\right)^{10} = \left(\frac{5}{6}\right)^{10} \approx 0.1615
\end{align*}

\paragraph*{Remark:}
The textbook uses $Y$, the number of faliures before the first success, instead of $X$, the number of trials until the first success.
This means that $Y = X-1$.

\subsubsection*{Hypergeometric Probability Distribution}

\paragraph*{Definition:}
Suppose you have a population of size $N$, in which each individual can be classified as either a success (S) or a failure(F).
The number of successes in the population is $M$. A sample of size $n$ is selected at random without replacement from the population.
Let $X$ be the number of successes in the sample. Then $X$ is a hypergeometric random variable with parameters $n$, $M$, and $N$:\\
\[
    X \sim \text{Hypergeometric}(n, M, N) \equiv X \sim HG(n, M, N)
\]

The pmf is denoted:
\[
    h(n,M,N)
\]

\paragraph*{Example:}
Identify the parameters $n$, $M$, and $N$ for the following situations:
\begin{itemize}
    \item The number of aces in a poker hand.
    \begin{itemize}
        \item $N=52$ (cards in deck)
        \item $M=4$ (aces in deck)
        \item $n=5$ (cards in hand)
    \end{itemize}
    \item The number of juniors in a 5-person committee chosen at random from a group of ten juniors and five seniors.
    \begin{itemize}
        \item $N=15$ (total students)
        \item $M=10$ (juniors)
        \item $n=5$ (committee members)
    \end{itemize}
    \item The number of green marbles you pull when randomly selecing 3 marbles from  a beg with 10 red,  8 green, and 2 blue marbles.
    \begin{itemize}
        \item $N=20$ (total marbles)
        \item $M=8$ (green marbles)
        \item $n=3$ (marbles selected)
    \end{itemize}
\end{itemize}

Let $X$ = the number of green marbles you pull when randomly selecting 3 marbles from a bag with 10 red, eight green, and two blue marbles.
Find the PMF of $X$.
\begin{align*}
    P(X=0) &= \frac{\binom{8}{0} \binom{12}{3}}{\binom{20}{3}} = \frac{220}{1140} = \frac{22}{114} \approx 0.19298 \\
    P(X=1) &= \frac{\binom{8}{1} \binom{12}{2}}{\binom{20}{3}} = \frac{8 \cdot 66}{1140} = \frac{528}{1140} = \frac{44}{95} \approx 0.46316 \\
    P(X=2) &= \frac{\binom{8}{2} \binom{12}{1}}{\binom{20}{3}} = \frac{28 \cdot 12}{1140} = \frac{336}{1140} = \frac{28}{95} \approx 0.29474 \\
    P(X=3) &= \frac{\binom{8}{3} \binom{12}{0}}{\binom{20}{3}} = \frac{56 \cdot 1}{1140} = \frac{56}{1140} = \frac{14}{285} \approx 0.04912
\end{align*}

\paragraph*{PMF:}
Generically, the pmf is given by:
\[
    h(x; n, M, N) = P(X=x) = \frac{\binom{M}{x} \binom{N-M}{n-x}}{\binom{N}{n}}, \quad \max\{n - (N-M), 0\} \leq x \leq \min\{n, M\}
\]

\begin{align*}
    E(X) &= n \left(\frac{M}{N}\right) \\
    V(X) &= n \left(\frac{M}{N}\right) \left(\frac{N-M}{N}\right) \left(\frac{N-n}{N-1}\right)
\end{align*}

\paragraph*{Example:}
Let $X$ = nuymber of aces in a 5 card poker hand. (a) Describe the distribution of X. (b) What is the probability that $x=2$?
(c) How many aces do you expect to have?

\paragraph*{a.}
\[
    X \sim HG(5, 4, 52)
\]

\paragraph*{b.}
\begin{align*}
    P(X=2) &= h(2; 5, 4, 52) \\
    &= \frac{\binom{4}{2} \binom{48}{3}}{\binom{52}{5}} \\
    &= \frac{6 \cdot 17296}{2598960} \\
    &= \frac{103776}{2598960} \approx 0.03993
\end{align*}

\paragraph*{c.}
\[
    E(X) = 5 \left(\frac{4}{52}\right) = \frac{5}{13} \approx 0.3846
\]

\subsubsection*{Relationship Between the Binomial and Hypergeometric Distributions}
\paragraph*{Remark:}
The hypergeometric distribution can be approximated by the binomial distribution when the population size $N$ and the 
number of successes $M$ is large relative to the sample size $n$.
